\section{관련 연구}\label{sec:related}

본 논문의 위치는 표~\ref{tab:positioning}이 정리한 네 갈래 가운데, \emph{추론 시점의 활성 스티어링}을 멀티-툴 선택 문제로 다시 가져오는 시도라는 점에서 읽어야 한다. 기존 연구는 대체로 ``도구 사용을 학습하느냐'', ``프롬프트로 유도하느냐'', ``추론 자원을 줄이느냐'', ``중간 표현을 직접 움직이느냐''로 나뉜다. 본 논문의 핵심은 이 가운데 마지막 갈래가 \emph{단일 개념 강조}에는 강했지만 \emph{순차적 멀티-툴 선택}에는 맞지 않았다는 점을 드러내고, 그 이유를 정지 연산자의 구조적 한계로 설명하는 데 있다.

첫 갈래인 tool-aware finetuning은 도구 사용 자체를 가중치에 학습한다. Gorilla~\citep{patil2023gorilla}는 대규모 API 호출 데이터로 function calling을 지도학습하고, ToolLLM~\citep{qin2024toolllm}는 실 API 집합 위의 tool-use trajectory를 통해 일반성을 넓히며, xLAM~\citep{zhang2024xlam}은 action model family로 확장한다. 이 계열은 강하지만 카탈로그가 바뀔 때마다 재학습 비용을 치른다. 둘째 갈래인 프롬프트 스캐폴딩은 학습 없이 reasoning chain과 행동 형식을 강제한다. ReAct~\citep{yao2023react}가 대표적이며, tool selection을 관찰-사고-행동 루프로 구조화한다. 그러나 프롬프트 길이와 호출 수가 task 복잡도에 따라 늘어나므로 추론 비용이 커진다. 셋째 갈래인 KV-cache 최적화는 또 다른 축이다. KIVI~\citep{liu2024kivi}와 같은 방법은 KV cache를 압축해 메모리와 latency를 줄이지만, \emph{어떤 도구를 고를 것인가}라는 의사결정 자체에는 개입하지 않는다. 따라서 본 논문의 질문과는 직교한다.

넷째 갈래인 활성 스티어링은 frozen model의 내부 표현을 추론 시점에 직접 수정한다. ITI~\citep{li2023iti}와 CAA~\citep{rimsky2024caa}는 residual stream 방향을 더하거나 빼고, PASTA~\citep{zhang2024pasta}는 post-softmax attention을 재가중하며, Focus Directions~\citep{zhu2025focusdirections}는 관련 context 쪽으로 주의를 유도한다. SEA~\citep{qiu2024sea}와 SEKA~\citep{li2026seka}는 spectral subspace를 이용해 activation 또는 key-space 개입을 구성한다. 이 계열의 공통점은 \emph{훈련 없이 행동을 바꾼다}는 데 있지만, 표준적인 배치 형태에서는 개입 자체가 대체로 \emph{정지적}이다. 같은 방향을 모든 디코딩 단계에 걸쳐 반복 주입하기 때문에, 이미 무엇을 선택했는지에 따라 개입이 바뀌지 않는다. prompt highlighting처럼 ``한 번 더 강하게 보라''는 문제에는 자연스럽지만, 멀티-툴 선택에서는 첫 도구를 계속 다시 밀어 주는 반복 편향으로 바뀔 가능성이 크다. 본 논문은 이 점을 \emph{tool-use setting에서 직접 측정한 결과}로 확인한다는 데서 기존 steering 논문과 갈라진다.

본 논문은 여기서 기존 활성 스티어링과 갈라진다. 첫째, 목표가 prompt highlighting이나 단일 개념 편집이 아니라 \emph{coverage와 비반복성}이 중요한 멀티-툴 selection이다. 둘째, 기저를 contrastive supervision으로 학습하지 않고 도구 카탈로그의 4-facet 구조에서 직접 구성한다. 셋째, 가장 중요한 차이로 개입을 전 레이어에 균일하게 두지 않고 초기 K와 전역 Q로 \emph{역할을 분리한 스케줄}을 사용한다. 즉 novelty는 ``더 좋은 projector''가 아니라 ``정지 연산자의 반복 편향을 피하기 위한 operator schedule''에 있다.

벤치마크 선택도 이 문제의식을 따른다. MetaTool Subtask4~\citep{MetaTool}는 각 질의에 정확히 두 개의 정답 도구가 있어 반복과 coverage를 가장 직접적으로 측정한다. $\tau^2$-bench~\citep{tau2bench}의 retail, telecom, airline은 정답 도구 수와 난도가 달라 short-horizon과 long-horizon을 함께 드러낸다. BFCL~\citep{BFCL}은 교차 검증용 function-calling 벤치마크다. 본문은 Qwen2.5-7B를 주 실험 모델로 두고, cross-model robustness는 Llama-3.1-8B에서만 제한적으로 다룬다. 이 구성은 ``멀티-툴에서 정지 스티어링이 왜 깨지는가''와 ``레이어별 역할 분리가 그 문제를 실제로 푸는가''를 가장 직접적으로 보여 주기 위한 것이다.
